\section{Derivations and proofs}
This appendix collects the mathematical backbone of the five tools in one
place. Each derivation is stated with its assumptions; where a step is an
approximation rather than a theorem, we say so.

\subsection{D1. Nearest-neighbor duplex free energy (Tool 6)}
\paragraph{Model.} For two RNA strands annealed in a given alignment, the
SantaLucia nearest-neighbor model assigns the standard free energy at
$37^{\circ}$C as
\begin{equation}
\Delta G^{\circ}_{37} \;=\; \Delta G^{\circ}_{\text{init}}
+ \sum_{k=1}^{m-1} \Delta G^{\circ}_{37}\!\left(\text{stack}_k\right)
+ \Delta G^{\circ}_{\text{term}},
\label{eq:nn}
\end{equation}
where $\text{stack}_k$ is the $k$-th adjacent base-pair doublet and the
per-doublet parameters are the experimentally fitted SantaLucia values.
\paragraph{Additivity.} Equation~\eqref{eq:nn} is an assumption, not a
theorem: it posits that the energy of a helix decomposes into
position-local contributions depending only on adjacent pairs. Its
empirical justification is the fit quality on optical-melting data; its
practical consequence is that duplex energy is computable in $O(m)$ per
alignment, so our screener can evaluate all $O(L^2)$ ungapped alignments
of two 20-nt spacers exactly, and delegate gapped/bulge structures to
ViennaRNA's \texttt{RNAcofold}, which minimizes the same objective over a
larger structure ensemble.
\paragraph{Screening statistic.} For guides $g_i, g_j$ we use
\begin{equation}
E_{ij} \;=\; \min_{a \in \mathcal{A}(g_i, g_j)} \Delta G^{\circ}_{37}(a),
\qquad
\text{flag}(i,j) = \mathbf{1}\!\left[E_{ij} \le \theta\right]
\;\vee\; \mathbf{1}\!\left[\text{shared PAM-proximal seed}\right],
\label{eq:flag}
\end{equation}
with $\theta = -8$~kcal/mol chosen so that the flag rate on random
spacers is near $10\%$ (measured $10.8\%$;
\texttt{results/realguide\_risk.json}). The min over alignments makes the
statistic conservative: any single strong annealing suffices to flag,
which is the correct bias for a screening (not a scoring) tool.

\subsection{D2. Two-proportion $z$-test for the same-gene enrichment (Tool 6 discovery)}
\paragraph{Setup.} Let $X_1 \sim \text{Binomial}(n_1, p_1)$ count flagged
same-gene pairs and $X_2 \sim \text{Binomial}(n_2, p_2)$ flagged
cross-gene pairs, $n_1 = n_2 = 30{,}000$. Under $H_0: p_1 = p_2$ the
pooled estimator is $\hat{p} = (X_1 + X_2)/(n_1 + n_2)$ and
\begin{equation}
Z \;=\; \frac{\hat{p}_1 - \hat{p}_2}
{\sqrt{\hat{p}(1-\hat{p})\left(\tfrac{1}{n_1} + \tfrac{1}{n_2}\right)}}
\;\xrightarrow{d}\; \mathcal{N}(0,1),
\label{eq:ztest}
\end{equation}
by the central limit theorem applied to the independent binomial counts.
\paragraph{Computation.} With $\hat{p}_1 = 0.22497$, $\hat{p}_2 = 0.15870$
(\texttt{results/realguide\_risk.json}): $\hat{p} = 0.19183$,
$\hat{p}(1-\hat{p}) = 0.15503$, the standard error is
$\sqrt{0.15503 \cdot 6.6667\mathrm{e}{-5}} = 3.215\mathrm{e}{-3}$, and
\begin{equation}
Z \;=\; \frac{0.22497 - 0.15870}{3.215\mathrm{e}{-3}} \;\approx\; 20.6,
\qquad p < 10^{-90}.
\end{equation}
\paragraph{Validity conditions.} The test treats sampled pairs as
independent. Pairs sharing a guide are weakly dependent, so the nominal
$p$-value is optimistic; the effect itself ($+6.6$ percentage points over
cross-gene, $+11.7$ over the random-spacer baseline $0.10812$) is large
relative to any plausible dependence correction, and the direction is
mechanistically predicted (transcripts of the same gene share exon
sequence, hence more reverse-complement complementarity). We report the
$z$-statistic as descriptive of effect size, not as an exact tail
probability.

\subsection{D3. Spearman correlation and the gene-clustered paired bootstrap (Tools 7, 10)}
\paragraph{Statistic.} For $n$ paired observations $(x_i, y_i)$ with rank
vectors $R_i, S_i$ (average ranks under ties), Spearman's coefficient is
the Pearson correlation of the ranks:
\begin{equation}
\rho \;=\; \frac{\sum_i (R_i - \bar{R})(S_i - \bar{S})}
{\sqrt{\sum_i (R_i - \bar{R})^2}\sqrt{\sum_i (S_i - \bar{S})^2}}.
\end{equation}
All our benchmark numbers are \emph{mean per-gene} Spearman:
$\bar{\rho} = \tfrac{1}{G}\sum_{g=1}^{G} \rho_g$, which weights genes
equally regardless of guide count.
\paragraph{Why a clustered bootstrap.} Guides within a gene share sequence
context and assay batch, so guide-level resampling would understate the
variance of $\bar{\rho}_{\text{leader}} - \bar{\rho}_{\text{challenger}}$.
Resampling \emph{genes} with replacement and recomputing both models'
per-gene correlations on the same resample preserves the pairing and the
cluster structure. The reported $p$-value is the fraction of $B = 10{,}000$
resamples with
\begin{equation}
\bar{\rho}^{(b)}_{\text{leader}} - \bar{\rho}^{(b)}_{\text{challenger}} \;\le\; 0,
\end{equation}
i.e.\ a one-sided test of ``leader at least as good''. This is exactly the
quantity a benchmark claim needs; with $G = 9$ genes the test is honest
but low-powered, which is why we report $p = 0.146$ (v7) and $p = 0.9996$
(v6) as \emph{failures to beat the bar}, not as near-misses
(\texttt{results/beat\_azimuth\_v6.json}, \texttt{results/beat\_azimuth\_v7.json}).

\subsection{D4. Microhomology deletion model and frameshift accounting (Tool 7)}
\paragraph{Lindel-style rate model.} For a deletion of length $\ell$
removing positions $[s, s+\ell)$ with microhomology of length $m \ge 0$
between the retained flanks, the model assigns an unnormalized rate
\begin{equation}
r(\ell, m, s) \;=\; \exp\!\left(\beta_0 + \beta_1 m + \beta_2 \ell + \beta_3\, \mathbf{1}[\ell = 1]\right),
\qquad
\Pr(\ell, m, s) = \frac{r(\ell,m,s)}{\sum_{\ell', m', s'} r(\ell',m',s')},
\label{eq:lindel}
\end{equation}
where the sum runs over all deletions compatible with the 65-nt context.
The exponential form guarantees positivity; normalization turns rates into
a categorical distribution over outcome classes.
\paragraph{Frameshift.} Coding consequence follows from length modulo 3:
\begin{equation}
\Pr(\text{frameshift}) \;=\; 1 - \sum_{\ell \equiv 0 \;(\mathrm{mod}\; 3)} \Pr(\text{del of length } \ell)
\;-\; \Pr(\text{in-frame insertion}).
\end{equation}
Because $+1$ insertions dominate the insertion mass and
$1 \not\equiv 0 \pmod 3$, the +1-insertion rate and the frameshift rate
are strongly coupled --- which is why our head-to-head
(\texttt{results/lindel\_headtohead.json}) reports both targets: a model
can win one and lose the other, and ours does.
\paragraph{Expected deletion length.} Under \eqref{eq:lindel},
\begin{equation}
\mathbb{E}[\ell] \;=\; \frac{\sum \ell\, r(\ell,m,s)}{\sum r(\ell,m,s)},
\end{equation}
a one-pass quantity we use as a summary feature for the CNN/GBT models.
No closed form exists because the compatible-deletion set depends on the
sequence; the computation is $O(W^2)$ in the context width $W = 65$.

\subsection{D5. Out-of-fold stacking: what it estimates and why in-fold stacking leaks (Tool 10, v5 vs v6)}
\paragraph{Claim.} Let base learners $f_1, \dots, f_M$ be trained with
$K$-fold cross-validation, and let $\tilde{f}_m(x_i)$ denote the
prediction for $x_i$ from the fold \emph{not} containing $i$. A
meta-learner $g$ trained on $\{(\tilde{f}_1(x_i), \dots, \tilde{f}_M(x_i), y_i)\}$
sees feature--label pairs whose joint distribution matches deployment
(features from models that never saw the label). In-fold predictions
$\hat{f}_m(x_i)$ from models that \emph{did} train on $(x_i, y_i)$ are
biased toward $y_i$; the meta-learner then learns to trust whichever base
learner overfits most.
\paragraph{Quantification on our split.} The effect is not hypothetical:
the v5 stack with in-fold deep features scored $0.370$ and lost to the
leader with $p = 1.0$; the v6 leak-free OOF stack on identical
architectures and data scored $0.490$ ($p = 0.9996$). The $+0.12$ gap is a
direct measurement of stacking leakage on this dataset
(\texttt{results/beat\_azimuth\_v5.json}, \texttt{results/beat\_azimuth\_v6.json}).
\paragraph{Caveat.} OOF stacking removes the \emph{label} leak but not the
\emph{hyperparameter} leak: tuning base learners on the same training pool
still biases the meta-features mildly upward. We accepted this residual
bias because the benchmark question (``can deep features beat the
same-split leader?'') is conservative under it --- the challenger still
lost.

\subsection{D6. CFD as a product of per-mismatch penalties (Tool 10 off-target replication)}
\paragraph{Derivation.} Model cleavage odds for an off-target with
mismatch set $M$ as factorizing over independent contributions:
\begin{equation}
\text{odds}(\text{active} \mid M) \;=\; \text{odds}_0 \prod_{i \in M} \lambda(p_i, b_i),
\end{equation}
where $p_i$ is the guide position and $b_i$ the mismatch base identity.
Taking $\Pr = \text{odds}/(1+\text{odds})$ and approximating
$\Pr \approx \text{odds}$ for small probabilities gives the CFD scoring
rule
\begin{equation}
\text{CFD}(M) \;=\; \prod_{i \in M} w(p_i, b_i),
\end{equation}
with $w$ the empirically fitted penalty table of Doench et al.\ 2016. The
independence assumption is known to be wrong for adjacent mismatches
(epistasis in the seed), which is one reason CFD saturates on
multi-mismatch targets; our replication uses it only as a reference
ordering check (CFD-style $0.582$ vs Hsu-style $0.556$ mean LOSO AUC,
winning 5 of 7 studies --- the published ordering;
\texttt{results/offtarget\_scorers.json}).

\subsection{D7. Tool 8 risk score: composition and invariance (Tool 8)}
\paragraph{Definition.} With repeat count $|R_s|$, worst span $d$,
complexity deficit $c$, and homopolymer term $h$,
\begin{equation}
S \;=\; 0.45\, \phi(|R_s|/20) + 0.30\, \phi(d/500) + 0.15\, c + 0.10\, h,
\qquad \phi(x) = \min(x, 1).
\label{eq:tool8}
\end{equation}
\paragraph{Properties.} (i) \emph{Boundedness}: $0 \le S \le 1$ because
each term is capped at its weight. (ii) \emph{Monotonicity}: $S$ is
nondecreasing in $|R_s|$ and $d$, matching the mechanistic direction of
the Kosicki et al.\ observation (more spanning repeats, longer possible
deletions $\Rightarrow$ higher risk). (iii) \emph{Saturation}: $\phi$
makes the score insensitive to repeat counts beyond 20 and spans beyond
500~bp; this is deliberate, since beyond those scales every locus is
simply ``bad'' and discrimination is useless for guide ranking.
\paragraph{What $S$ is not.} The weights are hand-set from the qualitative
weight of evidence in \cite{kosicki2018, wen2022}, not fitted, because no
per-sequence large-deletion \emph{rate} dataset exists. We therefore use
$S$ only for ranking and flagging (LOW/MODERATE/HIGH at $0.25/0.5$), and
we make no calibrated-probability claim anywhere in the tool, the tests,
or this manuscript.

\subsection{D8. Cost-asymmetric screening threshold (Tool 9)}
\paragraph{Decision rule.} Let $h(x) = \Pr(\text{hazardous} \mid x)$ be a
backend's assessed probability for oligo $x$, $C_{FN}$ the cost of
accepting a hazardous sequence, $C_{FA}$ the cost of blocking a safe one.
The Bayes-optimal accept rule minimizes expected cost:
\begin{equation}
\text{accept}(x) \iff C_{FA}\,(1 - h(x)) \;>\; C_{FN}\, h(x)
\iff h(x) < \tau^{\star} = \frac{C_{FA}}{C_{FA} + C_{FN}}.
\label{eq:bayes}
\end{equation}
\paragraph{Fail-closed as a limit.} For synthesis screening the cost
asymmetry is extreme ($C_{FN} \gg C_{FA}$), so $\tau^{\star} \to 0$ and
the optimal policy approaches ``accept only on positive evidence of
safety''. The wrapper implements exactly that limit: a backend failure is
treated as $h(x) = 1$ (block), never as $h(x) = 0$. Equation
\eqref{eq:bayes} also shows why a screening tool must not average verdicts
across backends: the expected-cost objective is not linear in $h$, so
averaging probabilities changes the decision boundary in the unsafe
direction whenever one backend is miscalibrated low.

\subsection{D9. Duplex-enrichment baseline under random spacers (Tool 6)}
\paragraph{Baseline construction.} The random-spacer flag rate
$0.10812$ in \texttt{results/realguide\_risk.json} comes from dinucleotide
-shuffled spacers paired against the same transcript pool, which preserves
per-guide base composition while destroying biological sequence identity.
The enrichment ratios are therefore
\begin{equation}
\text{ER}_{\text{same}} = \frac{0.22497}{0.10812} = 2.08,
\qquad
\text{ER}_{\text{cross}} = \frac{0.15870}{0.10812} = 1.47,
\end{equation}
and the composition-controlled excess of same-gene over cross-gene is
$0.22497 - 0.15870 = 0.0663$ in absolute flag rate. Because the shuffle
preserves GC, the excess cannot be attributed to composition; the
remaining explanation is shared exonic sequence within genes. This is the
quantified, falsifiable form of the Tool 6 discovery: any future guide-set
design rule that equalizes within-gene duplex risk should bring
$\text{ER}_{\text{same}}$ toward 1, and the screener measures whether it
does.
